% 本文件由 tools/build_book.py 从 dist/reports/assertions.json 生成，请勿手改。
\begin{longtable}{@{}p{0.06\textwidth}p{0.32\textwidth}p{0.22\textwidth}p{0.16\textwidth}@{}}
\caption{章节 5.3 数值断言（生成时间 2026-10-19T16:08:50-0400）}\\
\toprule
类别 & 断言 & 理论值 & 实测值 \\
\midrule
\endfirsthead
\toprule
类别 & 断言 & 理论值 & 实测值 \\
\midrule
\endhead
\bottomrule
\endlastfoot

A & 离散伴随 vs 中心差分(小步长 / 大步长)，以及标量模型的离散/连续闭式对照 & eps=\allowbreak{}1e-6 时 < 1e-9；eps=\allowbreak{}0.1 时 >=\allowbreak{} 1e-2(\allowbreak{}否则容差无意义)；与离散闭式 < 1e-12(\allowbreak{}伴随对离散问题精确)；与连续闭式的间隙按 16 倍缩小(\allowbreak{}4 阶) & duffing: eps=\allowbreak{}1e-6 -> 4.23e-11，eps=\allowbreak{}0.1 -> 3.15e-01；duffing-near-ridge: eps=\allowbreak{}1e-6 -> 1.09e-10，eps=\allowbreak{}0.1 -> 4.57e-01；标量模型离散闭式差 6.98e-14，连续/\allowbreak{}离散间隙 4.34e-10 -> 2.70e-11 -> 1.70e-12(\allowbreak{}步长比 16.09,\allowbreak{} 15.91) \checkmark \\
A & V 形曲线两段斜率、实测最优 eps 与理论比值 & 最优点左侧(\allowbreak{}舍入)斜率 -1；右侧(\allowbreak{}截断)斜率 2(\allowbreak{}中心)/\allowbreak{}1(\allowbreak{}单侧)；eps* 在扫描内部且与 u\textasciicircum{}(\allowbreak{}1/\allowbreak{}3)、\allowbreak{}u\textasciicircum{}(\allowbreak{}1/\allowbreak{}2) 同量级 & 中心 eps*=\allowbreak{}1.000e-06(\allowbreak{}u\textasciicircum{}(\allowbreak{}1/\allowbreak{}3)=\allowbreak{}4.806e-06，比值 0.208)，左段(\allowbreak{}舍入)斜率 -0.980、\allowbreak{}右段(\allowbreak{}截断)斜率 1.993，谷底误差 4.23e-11；单侧 eps*=\allowbreak{}3.162e-09(\allowbreak{}u\textasciicircum{}(\allowbreak{}1/\allowbreak{}2)=\allowbreak{}1.054e-08，比值 0.300)，左段 -0.997、\allowbreak{}右段 0.987；eps=\allowbreak{}1e-16 处误差 5.69e-01 \checkmark \\
A & 连续伴随 vs 离散伴随：线性重合边界与非线性阶数 & 常数雅可比模型上 < 1e-12；非线性按格式阶(\allowbreak{}rk4=\allowbreak{}4、\allowbreak{}heun=\allowbreak{}2、\allowbreak{}euler=\allowbreak{}1) & 线性 k=\allowbreak{}2: 1.11e-16，k=\allowbreak{}8: 3.33e-16，k=\allowbreak{}32: 3.33e-16；非线性 rk4: p=\allowbreak{}4.159(\allowbreak{}N=\allowbreak{}64-512，残差 2.35e-02 -> 2.32e-07)；heun: p=\allowbreak{}2.002(\allowbreak{}N=\allowbreak{}128-1024，残差 4.40e-02 -> 1.54e-04)；euler: p=\allowbreak{}0.998(\allowbreak{}N=\allowbreak{}2048-16384，残差 2.85e-02 -> 1.81e-03) \checkmark \\
A & 全存轨迹 vs 分块检查点的梯度一致性与存储/重算账 & 梯度逐位相同(\allowbreak{}相对差 0)；存储降到 1/\allowbreak{}sqrt(\allowbreak{}N) 量级，重算涨到正演的 2 倍附近 & N=\allowbreak{}256: 梯度相对差 0.0e+00，存储 3590->674(\allowbreak{}0.188 倍)，重算残差求值 1025->2048(\allowbreak{}1.998 倍于正演)；N=\allowbreak{}1024: 梯度相对差 0.0e+00，存储 14342->2370(\allowbreak{}0.165 倍)，重算残差求值 4097->8192(\allowbreak{}2.000 倍于正演) \checkmark \\
C & 伴随 vs 有限差分的正演次数比(k=2/8/32)与伴随次数的 k 无关性 & 正演次数比 =\allowbreak{} 2k(\allowbreak{}4/\allowbreak{}16/\allowbreak{}64)；伴随的求值次数与存储量与 k 无关(\allowbreak{}只与 N、\allowbreak{}dim 有关) & k=\allowbreak{}2: 伴随 rhs=\allowbreak{}385 jac=\allowbreak{}768 正演=\allowbreak{}1，有限差分 rhs=\allowbreak{}1536 正演=\allowbreak{}4，比值 4；k=\allowbreak{}8: 伴随 rhs=\allowbreak{}385 jac=\allowbreak{}768 正演=\allowbreak{}1，有限差分 rhs=\allowbreak{}6144 正演=\allowbreak{}16，比值 16；k=\allowbreak{}32: 伴随 rhs=\allowbreak{}385 jac=\allowbreak{}768 正演=\allowbreak{}1，有限差分 rhs=\allowbreak{}24576 正演=\allowbreak{}64，比值 64 \checkmark \\
A & loglog\_\allowbreak{}slope：幂律斜率、剔除非正值的行为与三种退化输入的拒答 & 剔除非正/\allowbreak{}非有限点后有效点 >=\allowbreak{} 2 才给斜率(\allowbreak{}0 在 log 轴上无信息)，否则 None & 幂律 -> 2.000000000000；含一个 0 -> 0.1982(\allowbreak{}剩 4 点)；全 NaN /\allowbreak{} 只剩一点 /\allowbreak{} 掩码后一点 -> None \checkmark \\
A & 端点组合帧数与 JSON 里的 NaN/Infinity 计数 & NaN/\allowbreak{}Infinity 计数 =\allowbreak{} 0 & 77 帧，计数 0(\allowbreak{}gradient\_field: 15 数组 /\allowbreak{} 28 meta；influence\_field: 10 数组 /\allowbreak{} 33 meta；gradient\_check: 13 数组 /\allowbreak{} 36 meta；forward\_adjoint\_playback: 11 数组 /\allowbreak{} 19 meta) \checkmark \\
A & 四帧的关键 meta 与 numcore 直接重算的逐位一致性 & 帧只是 numcore 结果的搬运工：同参数下的数字必须逐位相同 & 梯度场 max\_rel\_error=\allowbreak{}4.150502e-10；影响场 width\_ratio\_clean\_max=\allowbreak{}1.671154e-04,\allowbreak{} semidiscrete=\allowbreak{}1.677665e-06；梯度检验 eps*=\allowbreak{}1.000e-06,\allowbreak{} 斜率 -0.9742/\allowbreak{}1.9966；播放 growth=\allowbreak{}0.246764,\allowbreak{} 间隙=\allowbreak{}1.281e-05 \checkmark \\
A & eps 扫描窗口挪到截断段后'最优 eps'的报 null 行为 & 端点最优 -> eps\_opt\_* =\allowbreak{} null 且 *\_interior =\allowbreak{} false；内部最优 -> 给出数值 & eps\_min=\allowbreak{}1e-4 时 central/\allowbreak{}forward 均为 None,\allowbreak{} interior=\allowbreak{}False/\allowbreak{}False；默认窗口 1.00e-06(\allowbreak{}interior=\allowbreak{}true) \checkmark \\
A & 四个可视化的必备对照量与占位掩码的成对性 & requiredAux 指向的数组必须存在；*\_valid 掩码必须与数据同行且含 0 & gradient\_field: requiredAux (\allowbreak{}'error',\allowbreak{}) -> ['rel\_error\_log10',\allowbreak{} 'rel\_error']；influence\_field: requiredAux (\allowbreak{}'error',\allowbreak{}) -> ['width\_error',\allowbreak{} 'center\_error',\allowbreak{} 'semidiscrete\_abs\_error']；gradient\_check: requiredAux (\allowbreak{}'error',\allowbreak{}) -> ['error\_central',\allowbreak{} 'error\_forward']；forward\_adjoint\_playback: requiredAux (\allowbreak{}'residual',\allowbreak{}) -> ['sensitivity\_residual'] \checkmark \\
A & 灵敏度方向夹角、梯度相对误差与两种'最重要参数'排序 & 夹角与相对误差由差分误差决定(\allowbreak{}< 1e-3 度 /\allowbreak{} 1e-8)；原始梯度与弹性的排序可以不同 & 夹角 8.538e-07 度，相对误差 4.23e-11；原始梯度 [-1.811 -3.183 -0.794] -> 第 1 个分量最大；弹性 [-6.073 -1.067 -0.133] -> 第 0 个分量最大 \checkmark \\
A & 梯度场逐点对拍(正确 eps 与粗 eps)与共振脊所在阻尼 & 正确 eps 下逐点 < 1e-6；eps=\allowbreak{}0.1 的粗差分 > 1e-3；响应峰在最小阻尼处 & 逐点最大 1.41e-10、\allowbreak{}中位 5.80e-11(\allowbreak{}9 个箭头)；粗差分最大 6.88e-01；目标函数最大 1.0188 出现在 p2=\allowbreak{}0.030 \checkmark \\
A & 热方程伴随：正确转置 vs 差分，以及错转置在中心/贴边观测下的可见度 & 正确实现 < 1e-9；错转置在贴边观测时梯度 > 1e-6、\allowbreak{}场 > 1e-2，在中心观测时 < 1e-7 & obs=\allowbreak{}0.50,\allowbreak{} r=\allowbreak{}0.04: 正确 vs 差分 4.17e-12，错转置梯度差 1.64e-11，场相对差 2.29e-06；obs=\allowbreak{}0.90,\allowbreak{} r=\allowbreak{}2.00: 正确 vs 差分 6.74e-14，错转置梯度差 1.64e-05，场相对差 5.20e-02 \checkmark \\
A & 影响场宽度律、质心守恒与 Dirichlet 壁污染的量级 & 宽度平方斜率 /\allowbreak{} 2*alpha =\allowbreak{} 1(\allowbreak{}1e-3 内)；质心 <=\allowbreak{} 1e-12；触壁后偏差涨一个量级 & 斜率比 0.999748(\allowbreak{}实测 0.099975 /\allowbreak{} 理论 0.100000)；质心漂移 1.67e-16；干净窗口偏差 1.67e-04；tau=\allowbreak{}T 处 9.96e-03(\allowbreak{}宽度 0.14851 vs 0.15000)；污染行 32/\allowbreak{}61 \checkmark \\
C & PDE 伴随的时间阶：内部节点 vs 含端点，以及端点层的相对大小 & 内部 2 阶(\allowbreak{}CN)、\allowbreak{}含端点 1 阶(\allowbreak{}转置后的 Dirichlet 端点行)；两者比值 > 10 & 内部误差 3.76e-05 -> 1.47e-07，拟合阶 2.000；含端点误差 3.09e-04 -> 1.98e-05，拟合阶 0.991；端点层比值 67.19(\allowbreak{}含端点 3.95e-05 /\allowbreak{} 内部 5.87e-07) \checkmark \\
A & 伴随反向放大率随阻尼过零换向(实测 vs exp(-p2*T/2)) & |mu(\allowbreak{}0)|/\allowbreak{}|mu(\allowbreak{}T)| =\allowbreak{} 正向压缩率的倒数 =\allowbreak{} exp(\allowbreak{}-p2*T/\allowbreak{}2)(\allowbreak{}线性算例有闭式) & p2=\allowbreak{}-0.10: 19.9840(\allowbreak{}理论 20.0855)；p2=\allowbreak{}-0.05: 4.4791(\allowbreak{}理论 4.4817)；p2=\allowbreak{}+0.00: 0.9988(\allowbreak{}理论 1.0000)；p2=\allowbreak{}+0.05: 0.2217(\allowbreak{}理论 0.2231)；p2=\allowbreak{}+0.10: 0.0493(\allowbreak{}理论 0.0498)；最大相对差 9.98e-03 \checkmark \\
A & 本章 theta 格式 vs 4.1 的 theta\_\allowbreak{}step，以及 PDE 伴随梯度与差分的对照 & 同口径实现逐步一致(\allowbreak{}< 1e-12)；梯度与中心差分一致(\allowbreak{}< 1e-9) & 逐步最大偏差 9.71e-17；梯度 [0.0077102 0.0077102]，相对误差 1.54e-12(\allowbreak{}差分 4 次正演) \checkmark \\
\end{longtable}
